\section{Discusión}\label{sec:discusion}

\subsection{Comparación con la literatura}

El modelo seleccionado rinde por debajo de las cifras publicadas sobre estos datos: un
ROC-AUC de 0,762 en el conjunto de prueba frente al 0,997 de \textcite{oprea2022} y al
liderato cercano a 0,86 de la competencia de Zindi. La diferencia se explica en buena parte
por lo que mide cada cifra. El control entrenado con el historial completo, con el mismo
algoritmo y las mismas precauciones de partición, alcanza en el conjunto de prueba un ROC-AUC
de 0,879 y una precisión del 60,1\,\% en el 1\,\% superior de la lista, una cifra cercana a
la del liderato. La verificación por mitades muestra de dónde sale esa ventaja: la
segunda mitad del historial, la que contiene lo ocurrido después de la detección, rinde más
del doble que la primera. Esa parte del desempeño no existiría al puntuar clientes que
todavía no han sido inspeccionados.

En el caso de \textcite{oprea2022} se suma una segunda fuente probable de optimismo, porque
las variables de ventana se calcularon sobre el conjunto completo antes de partirlo. Con los
datos disponibles no es posible separar cuánto de su 0,997 corresponde a cada mecanismo,
pero la combinación de un único corte aleatorio con variables calculadas sobre todos los
clientes impide tomar esa cifra como referencia.

El hallazgo coincide con el de \textcite{comapuig2022} en datos de otra empresa y otro país:
un clasificador entrenado con etiquetas de inspección termina aprendiendo el proceso de
inspección. Los datos de STEG no traen un historial explícito de visitas; aquí la huella
apareció en cambios administrativos del contador y del contrato, que se registran como
cualquier otra factura. La consecuencia práctica es la misma: sin un control específico,
la validación cruzada premia al modelo que mejor reconoce a los clientes ya inspeccionados.

\subsection{Qué aprendió el modelo}

La variable con mayor contribución es la proporción de facturas con código de contador 207.
Es también la que más dudas plantea, porque ese código gana presencia en la segunda mitad del
historial de los clientes con fraude y se asocia con equipos más recientes. Tras la censura,
las variables del modelo todavía permiten predecir si el cliente tendrá un evento de equipo o
de contrato con un ROC-AUC de 0,871. Parte de la señal, por tanto, puede seguir anticipando
la inspección en lugar del fraude. El resultado en el grupo sin eventos acota ese riesgo: allí
no hay corte que delatar y el modelo mantiene un ROC-AUC de 0,758 en el conjunto de prueba,
por encima de la línea base de exposición.

El nivel y la dinámica del consumo suman el 37,8\,\% de la contribución. El percentil 90 del
consumo eléctrico, la pendiente anual y la proporción de caídas fuertes son variables que
una empresa puede auditar con su propio conocimiento del fraude, algo que no ocurre con un
modelo cuya señal principal fuera la antigüedad del cliente.

\subsection{Limitaciones}

La etiqueta registra fraude detectado y no tiene fecha. Un fraude cometido y nunca
inspeccionado figura como negativo, así que todas las métricas miden la capacidad de
anticipar qué inspección confirmará fraude y no la de detectar fraude cometido. Por la misma
razón no fue posible construir un corte temporal entre historial y etiqueta.

La partición de Zindi es aleatoria por cliente y no cronológica. El desempeño reportado no
mide la capacidad del modelo de generalizar a fraude futuro, que es lo que exigiría su uso en
una campaña real.

La censura es incompleta. Una inspección que confirma fraude sin cambiar el equipo ni el
contrato no deja evento que censurar. Entre las mitades del historial censurado persiste una
diferencia de 0,032 de PR-AUC (IC 95\,\% [0,022; 0,041]), compatible con la mayor relevancia de
lo reciente pero que no descarta un efecto residual de la inspección. La regla que decidió
qué variables con deriva de calendario se retiraban se fijó después de ver una versión
anterior de los resultados; por eso se evaluó el modelo sin el grupo B como control, y su
PR-AUC de 0,174 en el conjunto de prueba sigue por encima de la línea base de exposición.

Los datos son de Túnez. Los códigos de tarifa y de contador no tienen diccionario y el
consumo no tiene unidad, de modo que las variables más importantes no pueden interpretarse
en términos físicos ni trasladarse sin más a otra empresa. Aplicar el método en Colombia
requiere datos de las comercializadoras con la fecha y el resultado de cada inspección.

No se dispuso de costos de inspección ni del valor de la energía recuperada, así que no se
fijó un punto de operación ni se calculó el beneficio neto de cada tamaño de lista. La curva
de la Tabla~\ref{tab:curva-k} es el insumo para ese análisis, que queda pendiente de los
parámetros de costo de la empresa.
